\documentclass{article}


\title{MORE ADVANCED GIT}

\date{2024 DESAMBER 4}

\author{Ehsan Esmaeili}

\begin{document}
\maketitle
\newpage

\section*{More git commands}

\subsection*{Cloning}
\begin{enumerate}
    \item what the cloning means ? 

    cloning means you can save a copy of a project with all branch and commits , ets in your system as local .

    therefor you can add commits offline and send your new commit to orginal file of project later .
    \item what the commands make clone ?

    we can use this command to do it.
    
    git clone (url of project) 
    \item cloning a repasitury retain the commit history ?

    yes when you clone a ripasitury the history of repasitury are clone in your local project.
\end{enumerate}
\subsection*{Fork in open sourse project and git hub}
\begin{enumerate}
    \item means of fork :
    forking a project in github means you get a personal copy of project and make cheange in copy whisout changing in main project of anyone .

    this very helpfull  for open sourse project becaus encourages developer to help other in the open sourse project.
    
    \item process of forking :
    for forking you should go to page of repasitury at github and select fotk in page and fork will be save in your account .

    them you make change this fork and commit them , when you commit them they are will save in your account as your project.
    
    \item forking contibute in collaborate ptoject :
    \begin{itemize}
        \item encouraging expeimentation:
        developer can try any idea without worry about breaking the orginal project. 
        \item facilitating contributions:
        the developers can make pull request for suggestion to make project better of pull request will be accept the suggest will atomaticaly apply in project.

    \end{itemize}
    \item keeping update cheanges with orginal repasitury :
    we will at first fetch the main repasitury in local project them merg 
    the local changes in maked remote of main project and push remote to apply cheanges in main repasitury.
    \item contribute to the orginal repository after forking :
    to help project at first fork project so make changes and commit it ,
    them push commited changes in git hub.

    after we pushing commited chage we make a pull request and send it afer sending pull request and acces by manager of project the changes will be appled in repasitory.
    
\end{enumerate}
\subsection*{Git resetting}

\begin{enumerate}
    \item porpus of git reset command in git :
    it can help us when we want to cancelling new changes and commits and want to go to which commit is exist before.
    \item diffrece between git reset and git revert :
    git reset delete the commit which you want fix it and you return to before commit and rewrite the code to fix buges.

    git revert make a new commit that commit neutre commit which has buges .
    
    therefore better you use git revert in your projects .
    \item explaning the tree primary mode of git reset :

\begin{itemize}
    \item git reset --soft :

    this command use when you could return a commit but retain ready cheanges of commit .
    \item git reset --mixed (defult mode):this commit use when you could restage a commit but you could have this commit . 
    \item git reset --hard :
    this command use when you want completely discard all changes with out retain any things . 
    
\end{itemize}
    
\end{enumerate}

\section*{Merging and rebasing}

\subsection*{There are we have two link to see merging in project one and rebasing in project two . }

\begin{itemize}
\item link one :

\item link two :
\end{itemize}

\end{document}
